\section{The uniform cusp-weight mean on unitary period quotients}
\label{sec:sr-uniform-cusp-mean}

We prove the numerical mean estimate used by the covolume
construction.  The constants in this section are chosen before the
field, Hermitian space, lattice and arithmetic level.
The construction of an integral model with its tied period map is a
separate issue; no such model is inferred from the mean estimate.

Let \(E/F\) be a quadratic CM extension, put
\[
 [F:\mathbb Q]=f,\qquad [E:\mathbb Q]=d=2f,\qquad
 R=|\operatorname{disc}E|^{1/d},
\]
and fix an \(E\)-Hermitian space \(V\) of dimension \(N=p+q\),
of signature \((p,q)\) at every real place, where
\[
                   q\ge1,\qquad 3q+1\le N .
\]
Let \(\Lambda\subset V\) be a full \(\mathcal O_E\)-lattice.
We use the canonical homogeneous Hodge metric of a compatible
polarized Riemann-lattice realization.  Its alternating form \(B\) is integral on \(\Lambda\), its positive
Hermitian majorants have imaginary part \(B\), and
\[
 B(ax,y)=B(x,\overline a y)\quad(a\in\mathcal O_E).
\]
The \(\mathcal O_E\)-action is the one induced by the fixed
\(E\)-vector space.  In particular its operator norm, for every
compatible majorant, is at most
\(\max_\sigma|\sigma(a)|\).
Let \(\Gamma\) be any arithmetic subgroup of the integral
unitary group of this same lattice and form.
Write \(X_\Gamma=\Gamma\backslash G_\infty\), with invariant
probability measure, where
\(G_\infty=\prod_{\tau:F\hookrightarrow\mathbb R}U(p,q)\).
Passing to the quotient by a maximal compact subgroup makes no
difference to any function below.

For a majorant indexed by \(g\), let \(\|\cdot\|_g\) be its
Euclidean norm and
\[
                  \lambda(g)=
                     \min_{0\ne x\in\Lambda}\|x\|_g .
\]
Every minimum is positive and attained.  This follows from
discreteness of a full lattice and properness of a positive
Euclidean norm.  The minimum is a continuous function of \(g\):
on a relatively compact neighborhood only finitely many lattice
vectors can have norm below any fixed bound, and the minimum is
locally the minimum of those finitely many continuous functions.

For a nonzero saturated \(\mathcal O_E\)-submodule \(I\subset\Lambda\)
that is \(B\)-isotropic, let \(r=\operatorname{rank}_{\mathcal O_E}I\).
Its real rank is \(2fr\), and \(r\le q\).
Indeed isotropy on the \(\mathcal O_E\)-module implies
\(B(ax,y)=0\) for every \(a\in E\); nondegeneracy of the trace
pairing then gives isotropy for the underlying \(E\)-Hermitian
form.  Its dimension is at most the Witt index \(q\).

Use the normalized cusp covolume
\begin{equation}\label{eq:sr-cusp-covol-normalization}
                  n_g(I)=
                   \frac{\operatorname{covol}_g(I)}{R^{fr}},
 \qquad
 U(g)=\max\left(0,\sup_I-\log n_g(I)\right).
\end{equation}
This is the source's normalization, including its discriminant
factor.  Here is a direct check against its determinant definition.
Choose independent \(v_1,\dots,v_r\in I\) with
\(J=\sum_i\mathcal O_Ev_i\) of finite index in \(I\), and let
\(H_\sigma\) be the Gram matrix of their \(\sigma\)-eigenprojections.
If \(\sigma\) runs through one CM type, that definition is
\[
                     n_g(I)=
                   \frac{\prod_\sigma\det H_\sigma}{[I:J]} .
\]
For \(x,y\in I\), the imaginary part of their polarized inner
product is zero.  Applying this also to \(ax,by\) and using
nondegeneracy of the embedding matrix shows that the Gram matrices
for conjugate eigenprojections are complex conjugates.
They have the same positive determinant.
In a tower basis \(\omega_jv_i\), where \((\omega_j)\) is an
integral basis of \(E\), the real Gram matrix is obtained from the
block diagonal matrices \(H_\rho\), for all embeddings \(\rho\),
by the embedding matrix \((\rho(\omega_j))\).
Taking determinants therefore gives
\[
 \det\operatorname{Gram}_{\mathbb R}(J)
       =|\operatorname{disc}E|^r
                        \left(\prod_\sigma\det H_\sigma\right)^2 .
\]
Changing from \(J\) to \(I\) divides this determinant by
\([I:J]^2\).  Since
\(R^{2fr}=|\operatorname{disc}E|^r\), taking positive square roots
gives exactly \eqref{eq:sr-cusp-covol-normalization}.
There is no suppressed factor of two in this computation.
For example in \(E=\mathbb Q(i)\), take an isotropic real
unit vector whose two conjugate eigenprojections have squared
norm \(1/2\).  The real Gram determinant for its \(1,i\)
translates is \(1\), whereas the displayed expression is
\(4(1/2)^2=1\).  More generally the embedding determinant is
the determinant over \emph{all} complex embeddings; the two
conjugate projection blocks already contain these factors.
Thus the proof below concerns the actual determinant cusp weight
for the canonical tied Hodge metric.

\subsection{Short integral multipliers and isotropic vectors}

There are \(d\) independent algebraic integers
\(a_1,\dots,a_d\) satisfying
\begin{equation}\label{eq:sr-cusp-short-integers}
                 \max_\rho|\rho(a_i)|\le dR .
\end{equation}
For completeness use the Euclidean embedding whose squared norm is
\(\sum_\rho|\rho(x)|^2\).
The dual of the integer lattice is the conjugate embedding of the
codifferent.  If \(0\ne b\) belongs to the codifferent, ideal
containment gives
\(|N_{E/\mathbb Q}b|\ge|\operatorname{disc}E|^{-1}\).
The arithmetic--geometric mean inequality yields
\[
                  \sum_\rho|\rho(b)|^2\ge d/R^2 .
\]
The Euclidean transference inequality
\(\lambda_d(L)\lambda_1(L^*)\le d\)
therefore supplies \(d\) independent integral vectors of norm
at most \(\sqrt d\,R\), hence at most \(dR\).
We use only this standard transference inequality from
Banaszczyk, \emph{New bounds in some transference theorems in
the geometry of numbers}, Theorem~2.1,
\url{https://doi.org/10.1007/BF01445125}.
Coordinate norms are bounded by Euclidean norm, proving
\eqref{eq:sr-cusp-short-integers}.

Set
\[
                 v_0=2fR=dR,\qquad r_0=v_0^{-1/2},
 \qquad a=2fq,\qquad k=2f(N-1).
\]
Then \(k\ge3a\).
Any two lattice vectors of norm less than \(r_0\) are
\(E\)-orthogonal.  In fact
\[
 |B(a_i x,y)|\le
              \|a_i x\|_g\|y\|_g
          \le v_0\|x\|_g\|y\|_g<1 .
\]
The left side is the absolute value of an integer, so it is zero.
The \(a_i\) span \(E\) over \(\mathbb Q\); the same vanishing holds
for every element of \(E\), and the nondegenerate trace pairing
identifies this with \(E\)-Hermitian orthogonality.
Consequently all vectors of norm below \(t\le r_0\) lie in one
real space of dimension at most \(a\).
Disjoint balls of radius \(\lambda(g)/2\) in that space give
\begin{equation}\label{eq:sr-cusp-packing}
 \#\{0\ne x\in\Lambda:\|x\|_g<t\}
       \le(1+2t/\lambda(g))^a
       \le(3t/\lambda(g))^a
                   \quad(\lambda(g)<t\le r_0).
\end{equation}

A bounded divisibility condition will make the counting function
unfold into a finite sum of orbital integrals.
For \(0\ne x\in\Lambda\), define
\[
                \iota(x)=
          [\,\Lambda\cap Ex:\mathcal O_Ex\,].
\]
For a shortest vector \(x\), this index is at most
\[
                    C_0=(v_0\sqrt d)^d .
\]
To see this, the real \(d\)-dimensional lattice
\(L_x=\Lambda\cap Ex\) has minimum \(\lambda(g)\).
Minkowski's convex-body theorem gives
\(\operatorname{covol}(L_x)\ge
(\lambda(g)/\sqrt d)^d\); the weaker constant used here follows
by putting a cube of volume \(2^d\) inside the Euclidean ball
of radius \(\sqrt d\).
The subgroup generated by \(a_1x,\dots,a_dx\) has covolume
at most \((v_0\lambda(g))^d\), by Hadamard's determinant bound.
It is contained in \(\mathcal O_Ex\), so division of the two
covolumes proves the stated index bound.

Let \(\mathcal V\) consist of the nonzero \(E\)-isotropic
\(x\in\Lambda\) with \(\iota(x)\le C_0\), and define
\[
                    c_t(g)=
                  \#\{x\in\mathcal V:\|x\|_g<t\}.
\]
This is a nonnegative measurable function.
It is zero when \(t\le\lambda(g)\), is bounded by
\eqref{eq:sr-cusp-packing} when \(t\le r_0\), and is at least
one when \(\lambda(g)<t\le r_0\): a shortest vector then
belongs to \(\mathcal V\).

\subsection{The exact homogeneous mean}

We now prove
\begin{equation}\label{eq:sr-cusp-count-homogeneous}
                 \int_{X_\Gamma}c_t\,d\mu=A t^k
                       \quad(t>0)
\end{equation}
with some finite \(A\ge0\); no uniform bound on \(A\) is assumed.

The two reduction-theoretic background facts needed here are
finite covolume for arithmetic groups with no rational
characters, and finiteness of rational parabolic double cosets.
They are Borel--Harish-Chandra's finite-volume theorem and
Borel, \emph{Introduction aux groupes arithm\'etiques},
Theorem~15.5 and Proposition~15.6, applied by restriction of
scalars to \(\mathbb Q\).
We spell out their applications, since neither a uniform
covolume nor a quantitative cusp count is needed.

There are finitely many \(\Gamma\)-orbits of rational isotropic
\(E\)-lines: their stabilizers are the line parabolics.
On a fixed such line, \(I_\ell=\Lambda\cap\ell\) is a
fractional ideal times a fixed vector.
The condition \([I_\ell:\mathcal O_Ex]\le C_0\) leaves finitely
many possibilities for the principal submodule
\(\mathcal O_Ex\): its quotient is an integral ideal of bounded
norm after identifying the fixed ideal class.
Generators of each principal submodule form one
\(\mathcal O_E^\times\)-orbit.
A subgroup of finite index in these units occurs in the line
stabilizer inside \(\Gamma\).
Indeed choose a rational hyperbolic complement and write its
Levi action as
\(\operatorname{diag}(u,1,\overline u^{-1})\).
Units in a suitable congruence subgroup preserve the fixed
lattice.  Their intersection with \(\Gamma\) has finite index
by arithmetic commensurability with the integral stabilizer.
This argument does not require \(\Gamma\) itself to be a
congruence subgroup.
There are therefore only finitely many
\(\Gamma\)-orbits in \(\mathcal V\).

Fix an orbit representative \(x\).
Its real stabilizer has unipotent radical and reductive
quotient \(U(x^\perp/Ex)\).
The latter has anisotropic center over \(F\), and hence no
nontrivial \(F\)-rational characters; its arithmetic
quotient has finite volume.
The unipotent arithmetic quotient is compact.
Thus the stabilizer quotient occurring in the orbit
unfolding has finite volume.

At one infinite place the orbit of a nonzero isotropic
vector is the cone \(Q(z)=h(z,z)=0\), with its vertex deleted.
The unitary group is transitive on it, by extending the
vector to a hyperbolic pair.
Its invariant measure is a positive constant times
\(\delta(Q(z))\,dz\), equivalently surface measure divided
by \(\|\nabla Q\|\).
The group preserves Lebesgue measure and \(Q\), so this
measure is invariant.
Under \(z\mapsto tz\), \(dz\) contributes \(t^{2N}\)
and the delta factor contributes \(t^{-2}\).
It therefore scales by \(t^{2N-2}\).
The measure of its intersection with a bounded majorant
ball is finite: in polar coordinates its radial density
is \(r^{2N-3}dr\), integrable at zero, and the angular
section is compact and smooth.
Taking products over the \(f\) real places gives degree
\(2f(N-1)=k\).

Tonelli's theorem unfolds the nonnegative count of this
orbit into the stabilizer volume times its invariant
cone integral.
The preceding paragraph shows that the result is a finite
constant times \(t^k\).
Summing over the finitely many orbits, and dividing by
the finite positive volume of \(X_\Gamma\), proves
\eqref{eq:sr-cusp-count-homogeneous}.
If there are no rational isotropic lines, the count is
zero and the same formula holds with \(A=0\).

\subsection{A uniform logarithmic tail}

The unknown constant \(A\) disappears in a short argument.
Put \(Z=r_0/\lambda\) and \(X=A r_0^k\).
From the lower bound on \(c_t\) we have
\[
             \Pr(Z>u)\le\min(1,Xu^{-k})\quad(u\ge1).
\]
If \(X=0\), then \(Z\le1\) almost everywhere.
If \(0<X\le1\), integration of this tail gives
\(\mathbb E\log^+ Z\le X/k\le1/k\).
If \(X>1\), the packing upper bound gives
\[
 X=\mathbb E c_{r_0}
       \le3^a\mathbb E(Z^a{\bf1}_{Z>1})
       \le3^a\frac{k}{k-a}X^{a/k}.
\]
For the last inequality split the layer-cake integral
at \(u=X^{1/k}\).
Taking logarithms and integrating the same tail once
more gives
\[
 \mathbb E\log^+ Z
 \le\frac{\log X}{k}+\frac1k
 \le \frac{a\log3+\log(k/(k-a))}{k-a}+\frac1k
 \le3 .
\]
The last bound uses \(k\ge3a\), \(a\ge1\),
\(\log3<2\), and \(\log(3/2)<1\).
In particular all logarithmic quantities just integrated
are integrable.  We have proved, uniformly in the field,
lattice and level,
\begin{equation}\label{eq:sr-cusp-minimum-mean}
             \mathbb E\log^+(r_0/\lambda)\le3 .
\end{equation}

Minkowski applied to a cusp sublattice \(I\) of rank
\(s=2fr\) yields
\(\operatorname{covol}_g(I)\ge
(\lambda(g)/\sqrt s)^s\).
Using \eqref{eq:sr-cusp-covol-normalization} and \(r\le q\),
\[
 \begin{split}
 -\log n_g(I)
   &\le fr\{2\log(1/\lambda(g))+\log(2frR)\}\\
   &\le fq\{2\log^+(1/\lambda(g))+\log(2fqR)\}.
 \end{split}
\]
The latter expression is nonnegative and independent of \(I\).
It bounds the supremum defining \(U\), and proves that
this supremum is finite.
It is measurable, being a supremum over a countable
collection of submodules of a countable lattice.
Since \(r_0=v_0^{-1/2}\) and \(v_0\ge1\),
\[
          \log^+(1/\lambda)
             \le\log^+(r_0/\lambda)+\tfrac12\log v_0 .
\]
Consequently
\begin{equation}\label{eq:sr-cusp-mean-explicit}
         0\le\mathbb E U
           \le fq\{6+\log(2fR)+\log(2fqR)\}.
\end{equation}

For the principal hosts \(R\ge2\).
If \(B_0\ge4fN^2\), then \(B_0\ge64\) and \(2fq\le B_0\).
Thus both logarithms in
\eqref{eq:sr-cusp-mean-explicit} are bounded by
\(\log(B_0R)\), while \(\log(B_0R)\ge\log128>3\).
For any real \(L\) this proves the precise estimate
\begin{equation}\label{eq:sr-cusp-mean-four}
       \mathbb E U
           \le4fq\{\log(B_0R)+|L|\}.
\end{equation}
One can in particular take the source's \(L=\operatorname{hostLval}\).
The constant is exactly \(C_u=4\), chosen before the host.

\subsection{Passage to the tied Hodge integral}

We identify the measure used above with the geometric
integral in the covolume construction.
In a matrix ball chart \(Z\in M_{p,q}(\mathbb C)\), the
Hodge determinant metric has potential, up to a fixed
positive scale and an additive constant,
\[
                         -\log\det(1-Z^*Z).
\]
Its Hessian sends a tangent matrix \(T\) to the positive
form with left and right factors
\((1-ZZ^*)^{-1}\) and \((1-Z^*Z)^{-1}\).
The determinant of this Hessian is
\[
 \det(1-ZZ^*)^{-q}\det(1-Z^*Z)^{-p}
                   =\det(1-Z^*Z)^{-N}.
\]
Taking products over the CM type shows that the top
Hodge curvature measure is a fixed positive scalar
times the invariant unitary measure.
The scalar cancels in a normalized mean.

This also handles the source's finite collections of
period sheets and generic multiplicities.
The integral of \(U\) is a sum of invariant integrals
with positive sheet multiplicities; its denominator
is the same sum with \(U\) replaced by \(1\).
Dividing gives a positive weighted average of the
means already bounded in
\eqref{eq:sr-cusp-mean-four}.
Branch loci and toroidal boundary divisors have zero
top-dimensional measure.  The established integrability
of \(U\) permits their removal without changing the
integral.
A compatible isometry of Riemann lattices preserves
the submodules, Gram determinants, indices and
minimum appearing in
\eqref{eq:sr-cusp-covol-normalization}; thus the cusp
weight on every sheet is the same actual function
pulled back from its lattice fiber.

For any tied canonical Hodge realization of the
principal packet, with \(q=j'\), \(N=j+j'\) and
\(f=|\mathrm{CMType}|\), we therefore obtain
\[
 \frac{\displaystyle
          \int U\,c_1(\overline{\mathcal M})^{\dim X}}
      {\displaystyle
          \int c_1(\overline{\mathcal M})^{\dim X}}
       \le4f j'
          \{\log(B_0\operatorname{hostRd})
                    +|\operatorname{hostLval}|\}.
\]
This proves the analytic numerical assertion
\srcid{MeanU} for the canonical determinant cusp weight.
It supplies the uniform mean part of the covolume
producer.  It does not supply the remaining
\srcid{HtU} lower bound at every algebraic point, nor
the simultaneous integral root-tower and Hodge
comparison witnesses needed to build
\srcid{CovolumeBoundOfX}.
